\section{The two-stage factorisation of $\Phi$: native holomorphic
factorisation algebra versus specialised chiral shadow}
\label{sec:two-stage-functor}

\subsection{The conflation that the native/shadow split resolves}

The working table of Vol~III presents, for each $d$, a single
$E_n$-chiral-algebra structure as the \emph{output} of
$\Phi_d: \mathrm{CY}_d\text{-cat}\to\mathrm{ChirAlg}$: the abelian
Heisenberg at $d=1$, the $E_2$-braided Mukai at $d=2$, the $E_1$-ordered
chiral at $d=3$. The Gerstenhaber-bracket-degree argument,
$\deg[-,-]_G = 1-d$ on $\mathrm{HH}^{\bullet}(\mathcal{C})$, is a genuine
theorem: it is what lets one read off $E_\infty$, $E_2$, $E_1$ in the
three cases. CY-A$_3$ on a reference elliptic curve is a genuine
chiral-algebra-on-a-curve theorem.

The labelling is backwards. What the table calls ``native output'' is
the value on a small disc, or on a reference curve, of a canonical
object that lives one dimension up. That canonical object is the
holomorphic $E_d$-factorisation algebra $\Phi^{\mathrm{FA}}_d(X)$ on
$X$ itself --- Costello--Gwilliam locality plus Kontsevich--Tamarkin
$E_d$-formality; Costello--Li for holomorphic Chern--Simons on CY$_3$.
The $E_n$-chiral structure read off the table is what one obtains
\emph{after} specialising $\Phi^{\mathrm{FA}}_d$ along a
$(d-1)$-dimensional closed sub-object $\Sigma_{d-1}\subset X$ with
reference curve $C\subset X$, via factorisation homology
$\int_{\Sigma_{d-1}}\big|_{C}$. The shadow depends on the pair
$(\Sigma_{d-1},C)$; the native holomorphic factorisation algebra does
not.

\subsection{The two-stage factorisation}

\begin{definition}[Two-stage CY-to-chiral functor]
\label{def:two-stage-phi}
For $d\geq 1$ and a Calabi--Yau manifold $X$ of complex dimension $d$
with canonical holomorphic volume form $\Omega_X$, set
\[
\Phi_d \;=\; \mathrm{Sp}_{\Sigma_{d-1},C}\;\circ\;\Phi^{\mathrm{FA}}_d
\]
where the two stages are:

\emph{Stage~1 (canonical).} $\Phi^{\mathrm{FA}}_d:
\mathrm{CY}_d\text{-cat}(X) \longrightarrow
E_d\text{-}\mathrm{HolFA}(X)$. The source is the category of
perfect complexes on $X$ with its Serre-functor CY structure; the target
is $E_d$ holomorphic factorisation algebras on $X$ in the sense of
Costello--Gwilliam (Vol.~2). The functor exists and is pinned up to
contractible choice by the combination of Kontsevich--Tamarkin
$E_d$-formality of polydifferential operators (acting on the
holomorphic-factorisation-algebra avatar of Hochschild cochains) and
Costello--Gwilliam locality. At $d=3$ on $X=\mathbb{C}^3$ the
holomorphic Chern--Simons construction of Costello--Li provides the
explicit pre-factorisation algebra; at $d=1$ on $X=E_\tau$ the
Heisenberg vertex operator algebra is the factorisation algebra of
observables of free chiral bosons on the curve.

\emph{Stage~2 (choice-dependent).} $\mathrm{Sp}_{\Sigma_{d-1},C}:
E_d\text{-}\mathrm{HolFA}(X) \longrightarrow E_1\text{-}\mathrm{ChirAlg}(C)$.
Given a closed $(d-1)$-dimensional analytic cycle
$\Sigma_{d-1}\subset X$ with reference curve $C\subset X$ meeting
$\Sigma_{d-1}$ transversely in a finite set, set
\[
\mathrm{Sp}_{\Sigma_{d-1},C}\bigl(\Phi^{\mathrm{FA}}_d(X)\bigr)
\;\coloneqq\;
\biggl(\int_{\Sigma_{d-1}} \Phi^{\mathrm{FA}}_d(X)\biggr)\Big|_{C}.
\]
The factorisation homology $\int_{\Sigma_{d-1}}$ is the
Lurie--Francis integral of an $E_d$-algebra over a $(d-1)$-manifold;
restriction to $C$ is pullback along $C\hookrightarrow X$. The output is
an $E_1$-chiral algebra (one complex direction remaining) on $C$. A
different pair $(\Sigma'_{d-1},C')$ in general produces a different
$E_1$-chiral algebra.
\end{definition}

\begin{remark}[Stage 1 is pinned, stage 2 is a moduli]
Stage~1 lives in the contractible choice-space of Kontsevich--Tamarkin
formality quasi-isomorphisms; as an isomorphism class of functors,
$\Phi^{\mathrm{FA}}_d$ is unique. Stage~2 is parametrised by a moduli
$(\Sigma_{d-1},C)\in \mathrm{Cyc}_{d-1}(X)\times\mathrm{Curves}(X)$ --- a
\emph{real} moduli with nontrivial topology, because different cycles
give non-isomorphic specialisations. The ``many BKMs from one
CY$_3$'' picture is the statement that $\mathrm{Cyc}_2(K3\times E)$
supports several distinguished cycles, each giving a sibling
Borcherds--Kac--Moody denominator at the chiral-algebra level.
\end{remark}

\subsection{Per-dimension dictionary}

\begin{theorem}[Per-dimension native/shadow dictionary]
\label{thm:per-d-native-shadow}
For each $d\in\{1,2,3,5\}$, the native output $\Phi^{\mathrm{FA}}_d(X)$
on the CY$_d$ ambient $X$ and a canonical specialisation cycle
$\Sigma_{d-1}$ are:
\begin{center}
\begin{tabular}{c|c|c|c}
$d$ & Native $\Phi^{\mathrm{FA}}_d(X)$ & Specialisation cycle $\Sigma_{d-1}$ &
$E_1$-chiral shadow on $C$\\
\hline
$1$ & $E_1$-hFA on $E_\tau$ & point (trivial) & abelian Heisenberg\\
$2$ & $E_2$-hFA on $K3$ & $S^1\subset K3$ Nakajima cycle & $E_2$-braided
$H_{\mathrm{Muk}}$\\
$3$ & $E_3$-hFA on $K3\times E$ (CS-Li hCS) & $K3$ fibre & $E_1$-ordered
$\mathrm{CY\text{-}A}_3$ on $E$\\
$3$ & $E_3$-hFA on $A\times E$ & $T^4\subset A\times E$ torus fibre &
sibling $E_1$-ordered on $E$\\
$5$ & $E_5$-hFA on $K3\times K3\times E$ & $K3\times K3$ & $E_1$-ordered
on $E$\\
\end{tabular}
\end{center}
The $E_2$-braided shadow at $d=2$ is the value at a $1$-dimensional
stratum of the $E_2$ factorisation algebra, recovered by the
Lurie--Francis--Kapranov theorem $\int_{S^1} E_2 \simeq E_1$ with the
$E_2$-braiding promoted to the monodromy data of an $E_1$-module
category on the Nakajima circle. The $E_1$-ordered shadow at $d=3$ is
$\int_{\Sigma_2}E_3\simeq E_1$ by Dunn--Lurie additivity
$E_3\simeq E_1\otimes E_2$ combined with
$\int_{\Sigma_2}(E_1\otimes E_2)=E_1\otimes \int_{\Sigma_2}E_2$ and the
vanishing of $\int_{\Sigma_2}E_2$ on a closed real surface.
\end{theorem}

\begin{proof}[Proof sketch]
Stage 1 is Kontsevich--Tamarkin. For stage 2 use Dunn--Lurie additivity
$E_{a+b}\simeq E_a\otimes E_b$ iteratively, together with
Lurie--Francis--Kapranov for factorisation homology of $E_n$-algebras
over disk-stratified manifolds. At $d=3$, $\Sigma_2=K3$, the
$E_3$-structure on $\Phi^{\mathrm{FA}}_3(K3\times E)$ writes locally as
$E_1\otimes E_2$ with the $E_1$ along $E$-direction and the $E_2$ along
the $K3$-direction; integrating the $E_2$ out over the closed $K3$
leaves the $E_1$ on $E$. The order-$1$ shadow depends on the $K3$
fibre class through the spectral parameter of the induced Yangian, but
the braided $R$-matrix structure of the $K3$-factor is what remembers
the $E_2$ Dunn summand --- this is where the Maulik--Okounkov stable
envelope enters as the explicit $R$-matrix of the $E_2$ component.
\end{proof}

\subsection{Sibling BKMs from a single $\Phi^{\mathrm{FA}}_3$}

\begin{corollary}[Many BKMs from one CY$_3$]
\label{cor:many-BKMs}
Let $X=K3\times E$ and $\Phi^{\mathrm{FA}}_3(X)$ its native $E_3$
holomorphic factorisation algebra. The following $E_1$-chiral
shadows all arise as $\mathrm{Sp}_{\Sigma_2,E}\bigl(\Phi^{\mathrm{FA}}_3(X)\bigr)$
for distinct choices of $\Sigma_2$:
\begin{enumerate}
\item $\Sigma_2=K3$ (fibre): the Gritsenko $\Delta_5$ Borcherds-lifted
$\mathfrak{g}_{\Delta_5}$ GKM superalgebra; $\kappa_{\mathrm{BKM}}=5$.
\item $\Sigma_2=E\times S^1_{\mathrm{sec}}$ for $S^1_{\mathrm{sec}}$ a
section circle of $E\to\mathrm{pt}$: a vertex-algebraic sector of the
Fake Monster at rescaled $\kappa_{\mathrm{BKM}}$ compatible with the
Borcherds--Harvey--Moore Igusa-$\Phi_{10}$ lift.
\item $\Sigma_2$ a Humbert surface in $\mathrm{Hilb}^2(K3)$ pulled back
along an isogeny: a sibling paramodular lift at weight $c_N(0)/2$ for
$N\in\{2,3,4,6\}$; $\kappa_{\mathrm{BKM}}\in\{2,1,1,1\}$.
\end{enumerate}
Each specialisation is a different theorem; there is not one $\Phi_3$
that outputs all of them, nor six applications of $\Phi_3$ to one
category, nor a conflation of $\Phi_3$ with the four invariants
$\{2,3,5,24\}$.
\end{corollary}

The statement sharpens three prior readings: (a) the ``six routes to
$G(K3\times E)$'' are six different specialisations of one
$\Phi^{\mathrm{FA}}_3$, not six applications of a functor $\Phi_3$; (b)
the $\{2,3,5,24\}$ spectrum arises from four different constructions
because the four constructions use different cycles $\Sigma_2$; (c) the
``Borcherds weight universality'' $\kappa_{\mathrm{BKM}}(\Phi_N)=c_N(0)/2$
is stable across $N$ precisely because the $N$-dependent content lives in
the choice of $\Sigma_2$ (a paramodular Humbert surface indexed by $N$),
not in the native $\Phi^{\mathrm{FA}}_3$.

\subsection{CY-A$_3$ is one specialisation, not a universal theorem}

\begin{remark}[Scope of CY-A$_3$]
\label{rem:cy-a3-specialisation}
The CY-A$_3$ equivalence, stated as a chiral-bar-cobar equivalence on
the elliptic reference curve $E$, is the specialisation
$\mathrm{Sp}_{K3,E}\bigl(\Phi^{\mathrm{FA}}_3(K3\times E)\bigr)$. It is a
theorem about a specific specialisation; it is not the statement that
$\Phi_3$ outputs a unique $E_1$-chiral algebra per category. Sibling
specialisations give different equivalences on different reference
curves. The native $\Phi^{\mathrm{FA}}_3$ is the uniform object; the
family of CY-A$_3$-type equivalences is indexed by the
$(\Sigma_2,C)$-moduli.
\end{remark}

\begin{remark}[Why the shadow labels were backwards]
The operadic channel ($E_n$-structure degree) and the cyclic/framing
channel ($(d-1)$-cycle choice) were collapsed onto a single column in
the original table. Separating them makes four invariants naturally
appear: $\kappa_{\mathrm{cat}}=\chi(\mathcal O_X)$ is a Stage~1 invariant
(property of $\Phi^{\mathrm{FA}}_d$, Künneth-multiplicative); $\kappa_{\mathrm{BKM}}$
is a Stage~2 invariant (property of the specific cycle, hence of
$c_N(0)/2$); $\kappa_{\mathrm{fibre}}$ is the difference between a
particular cycle and the ambient; $\kappa_{\mathrm{ch}}$ is the
Hodge-supertrace identification attached to the $E_1$-shadow.
\end{remark}

\subsection{Reconciliation with the working tables}

The revisions the two-stage reading forces on Vol~III are three:

\emph{(1) Rename the ``output'' column.} What the tables call
``output'' must be renamed ``$E_1$-shadow on $C$ via
$\mathrm{Sp}_{\Sigma_{d-1},C}$''. A second column must name the native
$\Phi^{\mathrm{FA}}_d(X)$.

\emph{(2) Make $(\Sigma_{d-1},C)$ explicit.} Every theorem asserting a
specific chiral-algebra output at $d\geq 2$ must name the cycle.
CY-A$_3$ on $E$ reads $\mathrm{Sp}_{K3,E}$ in the new notation; the
Borcherds--Harvey--Moore chapter reads $\mathrm{Sp}_{\Sigma_2^{\mathrm{Fake}},E}$
where the Fake Monster cycle is the Niemeier-class fixed locus.

\emph{(3) Replace the ``six routes'' paragraph.} Six routes to
$G(K3\times E)$ become six specialisations of one
$\Phi^{\mathrm{FA}}_3(K3\times E)$. The six cycles are catalogued:
$K3$-fibre, paramodular Humbert $N=2$, $\ldots$, each giving a sibling
BKM with $\kappa_{\mathrm{BKM}}$ reading off the corresponding
$c_N(0)/2$.

The cost of the rewrite is linear in the number of chapters asserting a
``CY-to-chiral output''; the dividend is that the multiplicity of
sibling BKMs stops being a coincidence and becomes a theorem about
$\mathrm{Cyc}_2$-moduli of $\Phi^{\mathrm{FA}}_3$.

\subsection{Claims and their status}

\begin{description}
\item[\ClaimStatusTheorem] Definition~\ref{def:two-stage-phi} and the
Stage~1 canonicality (Kontsevich--Tamarkin + Costello--Gwilliam).

\item[\ClaimStatusTheorem] Theorem~\ref{thm:per-d-native-shadow} at
$d=1,2$ with explicit references (Heisenberg; $H_{\mathrm{Muk}}$).

\item[\ClaimStatusDerivation] Theorem~\ref{thm:per-d-native-shadow} at
$d=3$ via Dunn--Lurie additivity. The sketch is rigorous; the full
proof requires the Costello--Li hCS pre-factorisation algebra and
its Stokes--Hochschild dualisation to the reference curve $E$.

\item[\ClaimStatusConjectured] The full catalogue in
Corollary~\ref{cor:many-BKMs} at points (2) and (3): the Fake-Monster
sibling and the paramodular-Humbert sibling are expected but require
the explicit identification of $\Sigma_2^{\mathrm{Fake}}$ inside
$K3\times E$ and of the Humbert surfaces as pullbacks along isogenies.

\item[\ClaimStatusDerivation] The scope corrections in
Remark~\ref{rem:cy-a3-specialisation}: CY-A$_3$ as written in
\texttt{chapters/theory/cy\_to\_chiral.tex} is an
$\mathrm{Sp}_{K3,E}$-specialisation statement.
\end{description}
